\chapter{Recomendaciones}
\label{cap:recomendaciones}

Las recomendaciones que siguen se desprenden de los resultados y de las limitaciones señaladas en las
Secciones~\ref{sec:limitaciones} y~\ref{sec:disc-limitaciones}.

\begin{itemize}
\item \textbf{Contraste de la hipótesis:} La hipótesis del trabajo (Capítulo~\ref{cap:introduccion}) se
planteó a partir de estos resultados, cuando ya se los conocía. Su examen muestra que es compatible con
ellos, pero esa compatibilidad no equivale a ponerla a prueba. Se recomienda contrastarla con datos
nuevos y con una regla de decisión fijada de antemano, que indique qué resultado la sostiene y cuál la
contradice.

\item \textbf{Confirmación del hallazgo en detección:} La ventaja de la propuesta sobre el Top-Hat
clásico fue mínima y dependió de la semilla. La tasa equilibrada de la propuesta, que promedia el acierto
sobre los objetos de cada cámara, fue 3.0 puntos mayor que la del Top-Hat clásico, con un intervalo de
confianza de 0.0 a 6.0. Con el \textit{bootstrap} salvaje por grupos, el intervalo fue de $-3.5$ a 5.5
e incluyó el cero. Según el par de semillas, la diferencia fue de $-3.1$ a 7.8 puntos. Confirmar esa
ventaja exige, ante todo, más grupos independientes en la prueba, de tamaño más parejo, armados por
secuencia de video y no solo por cuadros contiguos.
También conviene entrenar con más semillas por entrada, porque las tres diferencias negativas se dieron
frente a una misma semilla del Top-Hat clásico. Las luces, por su parte, deberían provenir de más tramos
de video: en la prueba, un solo grupo reunió casi dos de cada tres. Por último, la comparación debería
repetirse con otros detectores y con más épocas de entrenamiento. En nueve de los diez detectores
seguidos época por época, la mAP50 sobre los cuadros de seguimiento todavía subía en
las últimas cinco épocas. En esa misma medida, entre las épocas 30 y 40 cambió además el orden de la propuesta frente a la
pirámide de Laplace y la DWT (db4).

\item \textbf{Ajuste con un mismo criterio:} Los métodos comparativos se usaron con sus parámetros por
defecto, mientras que el radio de la propuesta se eligió sobre los mismos veinte pares con que se evaluó
la calidad. Se recomienda ajustar todos los métodos, la propuesta incluida, con un mismo criterio y
sobre pares que luego no entren en la evaluación. Sería útil, además, medir la calidad en otros
conjuntos, cuyas imágenes visibles provengan de cámaras comunes y no de cámaras con intensificador de
imagen.

\item \textbf{Aptitud del ajuste:} La grilla mostró que la aptitud, la función que el enjambre
busca maximizar, era mayor con pesos por debajo del rango de búsqueda: el mejor peso de cada radio
estuvo entre 0.05 y 0.07. Con pesos tan bajos, la regla de la Ecuación~\eqref{eq:regla-tophat} deja la
fusión cerca del promedio simple. Esa preferencia por pesos bajos podría deberse a que dos de los tres
términos de la aptitud (Ecuación~\ref{eq:aptitud}) premian el parecido con las fuentes. Se recomienda definir una
aptitud en la que ese parecido no domine, calcularla sobre pares ajenos a la evaluación y tomar el
promedio simple como referencia.

\item \textbf{Otros métodos comparativos:} Este trabajo se centró en un operador de una sola escala.
Como se indicó en la Sección~\ref{sec:limitaciones}, dejó fuera la fusión con redes neuronales y los
métodos morfológicos multiescala de \textcite{bai2013} y de \textcite{roman2020}. Incluirlos permitiría ubicar la propuesta
frente a esas dos familias.

\item \textbf{Revisión humana de los objetos de una sola cámara:} Un objeto se contó como visto por una
sola cámara si lo encontraban al menos dos de los tres detectores de esa cámara y a lo sumo uno de los
de la otra. Con ese criterio, un objeto que muestran ambas cámaras pudo contarse como de una sola cámara
si los detectores de la otra no lo encontraron. Conviene confirmar esa condición con una revisión
humana.

\item \textbf{Radio relativo al tamaño de la imagen:} El radio se expresó en píxeles, de modo que la
parte de la imagen que abarca depende de su resolución. Con $r = 25$, el disco de 51 píxeles de ancho
cubre entre el 8.9 y el 20.6~\% del lado menor de las imágenes del TNO. Si el radio se expresara en
proporción al tamaño de la imagen, abarcaría detalles de igual tamaño relativo en todas ellas.

\item \textbf{Fusión con la imagen visible en color:} La imagen visible se usó en escala de grises.
Queda por estudiar una fusión que combine el brillo de las dos imágenes y conserve el color de la
visible. Esa fusión interesa sobre todo para las luces, la clase en que la cámara común más aventaja a
la térmica.
\end{itemize}
